\textbf{Ejemplo 6}\\
Una industria tiene actualmente contratado un préstamo con una corporación financiera a la tasa del TCC + 3 puntos. ¿Cuál debe ser el spread en puntos básicos de forma tal que financieramente sea indiferente el préstamo en la corporación financiera o en el mercado de Londres? con estos índices:

\begin{itemize}
	\item $TCC = 15,3\% \textit{ nata}$
	\item $j_{dev} = 22\% \textit{ naav}$
\end{itemize}

Interés devaluación del peso Colombiano con respecto a la libra esterlina Libor

\begin{itemize}
	\item $Libor= 5,2\% \textit{ naav}$. Tasa de interés del mercado europeo
\end{itemize}

\textbf{Solución.}

\renewcommand{\arraystretch}{1.5} % Margenes de las celdas

\begin{longtable}[H]{|p{0.3750\linewidth}|p{0.3750\linewidth}|p{0.25\linewidth}|}
		%Creamos una linea horizontal
		\hline
		%Definimos el color de la primera fila
		\rowcolor[HTML]{FFB183}
		%%%%% INICIO ASIGNACIÓN FECHA FOCAL %%%%%%%
		%%%%%%%%%% INICIO TITULO
		%Lo que se hace aquí es mezclar las 3 columnas en una sola
		\multicolumn{3}{|c|}{\cellcolor[HTML]{FFB183}\textbf{1. Asignación período focal}}                                                                \\ \hline
		%%%%%%%%%% FIN TITULO
		\multicolumn{3}{|c|}{ $pf:1 \textit{ pav} \equiv 4 \textit{ ptv}$}
		\\ \hline
		
		%%%%% INICIO DECLARACIÓN DE VARIABLES %%%%%%%
		%%%%%%%%%% INICIO TITULO
		%Lo que se hace aquí es mezclar las 3 columnas en una sola
		\multicolumn{3}{|c|}{\cellcolor[HTML]{FFB183}\textbf{2. Declaración de variables}}                                                              \\ \hline
		%%%%%%%%%% FIN TITULO
		%%%%%%%%%% INICIO SUBTITULO
		\multicolumn{1}{|c|}{\cellcolor[HTML]{FFB183}\textbf{Valores monetarios}}    & 
		\multicolumn{1}{|c|}{\cellcolor[HTML]{FFB183}\textbf{Valores no monetarios}} &
		\multicolumn{1}{|c|}{\cellcolor[HTML]{FFB183}\textbf{Incógnitas}}  \\ \hline
		%%%%%%%%%% FIN SUBTITULO
		%%%%%%%%%% INICIO DE MATEMÁTICAS
		%Cada & hace referencia al paso de la siguiente columna
		$\textit{Para el presente ejercicio }  $\textit{no se utilizaron valores monetarios}$ $& $i_{dev} = 22\% \textit{ naav}$  & $X = ? \% \textit{ nata} $      \\
		& $TCC = 15,3\% \textit{ nata}$  &                                    \\
		& $libor = 5,2\% \textit{ naav}$ & \\ \hline
		
		%%%%%%%%%% FIN DE MATEMÁTICAS
		%%%%% FIN DECLARACIÓN DE VARIABLES
		
	\rowcolor[HTML]{FFB183}
	\multicolumn{3}{|c|}{\cellcolor[HTML]{FFB183}\textbf{3. Declaración de fórmulas}} \\ \hline
	\multicolumn{3}{|c|}{$i_{dev} + Libor + X = TCC + 3 \textit{ Ecuación de equivalencia}$} \\ \hline
	
	\rowcolor[HTML]{FFB183}
	\multicolumn{3}{|c|}{\cellcolor[HTML]{FFB183}\textbf{4. Desarrollo matemático}} \\ \hline
	\multicolumn{3}{|p{\textwidth}|}{
		\begin{center}
			$TCC + 3 puntos = 15,3\% \textit{ nata} + 3\% \textit{ nata} = 18,3\% \textit{ nata}$
			
			$j_{crédito} = TCC + 3 puntos = 15,3\% \textit{ nata} + 3\% \textit{ nata} = 18,3\% \textit{ nata}$
			
			$i = i_1 + i_2 + i_1 \cdot i_2 $ Equivalencia de tasas de referencia
			
			$j_c = 18,3\% \textit{ nata}$
			
			$j_c =? \textit{ naav}$ en Londres
			
			$18,3\% \textit{ nata} = 20,601\% \textit{ naav}$
		\end{center}
		
		Teniendo la tasa periódica vencida utilizamos la fórmula de las tasas equivalentes
		
		$(1 + i_1)^{m_1} = (1 + i_2)^{m_2}$, donde $m_1 = 1 \textit{ pav}$
		
		$(1 + i_1)^1 = (1 + 0,183)^4$
		
		$i_{dev} + Libor + X = 22\% + (5,2 + X)\%$
		
		Crédito equivalente = $22\% \textit{ naav} + (Libor + X)\% \textit{ naav} =? \% \textit{ naav}$
		
		Por la fórmula de combinación de tasas, por ser tasas naav
		
		$0,28344 + 0,0122 + X = 0,20601$ (Despejar la X)
		
		$X = -6,35\%$
	} \\ \hline
	
	\rowcolor[HTML]{FFB183}
	\multicolumn{3}{|c|}{\cellcolor[HTML]{FFB183}\textbf{5. Respuesta}} \\ \hline
	\multicolumn{3}{|p{\textwidth}|}{
		$X = -6,35\%$
	} \\ \hline
\end{longtable}